\documentclass[12pt,a4paper]{article}

% Detecção de motor e fontes
\usepackage{iftex}
\ifpdftex
  \usepackage[T1]{fontenc}
  \usepackage[utf8]{inputenc}
  \usepackage{lmodern}
\else
  \usepackage{fontspec}
\fi

% Idioma
\usepackage[brazilian]{babel}

% Geometria da página
\usepackage[top=2cm, bottom=2cm, left=2cm, right=2cm]{geometry}

% Pacotes Matemáticos
\usepackage{amsmath, amsfonts, amssymb}

% Cores e Caixas de texto
\usepackage{xcolor}
\usepackage[most]{tcolorbox}
\newtcolorbox{instrucoes}{
  colback=black!5,
  colframe=black,
  coltext=black,
  arc=2mm,
  boxrule=0.8pt,
  title=Instruções,
  fonttitle=\bfseries
}

% Listas e Questões
\usepackage{enumitem}
\newlist{questoes}{enumerate}{2}
\setlist[questoes,1]{label=\textbf{\arabic*.}, leftmargin=1.5em, itemsep=2em}
\setlist[questoes,2]{label=\alph*), leftmargin=1.5em}

% Cabeçalho Manual
\newcommand{\cabecalho}{
  \noindent
  \begin{minipage}[t]{0.5\textwidth}
    \vspace{0pt}
    \textbf{Instituição:} Nome da Escola/Universidade \\[4pt]
    \textbf{Professor:} Nome do Docente \\[4pt]
    \textbf{Disciplina:} Matemática
  \end{minipage}%
  \begin{minipage}[t]{0.5\textwidth}
    \vspace{0pt}
    \raggedleft
    \begin{tabular}{@{}ll@{}}
      \textbf{Data:}  & \_\_\_/\_\_\_/\_\_\_\_\_ \\[4pt] 
      \textbf{Turma:} & \_\_\_\_\_\_\_\_\_\_\_\_ \\[4pt]
      \textbf{Peso:}  & 10,0 \hspace{0.5cm} \textbf{Nota:}
    \end{tabular}
  \end{minipage}
  \vspace{0.5cm}
  \begin{center}
    \rule{\textwidth}{0.4pt} \\
    \vspace{0.2cm}
    {\large \textbf{AVALIAÇÃO}} \\
    \rule{\textwidth}{0.4pt}
  \end{center}
  \vspace{0.3cm}
  \noindent \textbf{Aluno(a):} \hrulefill
  \vspace{0.8cm}
}

\begin{document}

\cabecalho

\begin{instrucoes}
  \begin{itemize}[leftmargin=1.5em]
    \item As respostas finais das questões objetivas devem ser feitas a caneta (azul ou preta), enquanto as questões discursivas e os cálculos podem ser feitos a lápis.
    \item Todas as questões, inclusive as objetivas, exigem a apresentação do desenvolvimento. Respostas sem o respectivo desenvolvimento matemático serão desconsideradas ou sofrerão desconto na nota.
    \item O uso de calculadoras é permitido.
    \item Mantenha a organização e a legibilidade dos cálculos.
  \end{itemize}
\end{instrucoes}

\vspace{0.5cm}

\begin{questoes}

  % Questão de Múltipla Escolha
  \item (2,5 pontos) Enunciado de uma questão objetiva genérica de múltipla escolha.
  \begin{enumerate}[label=\alph*)]
    \item Primeira alternativa
    \item Segunda alternativa
    \item Terceira alternativa
    \item Quarta alternativa
    \item Quinta alternativa
  \end{enumerate}

  % Questão Aberta com Subitens
  \item (3,0 pontos) Enunciado de uma questão discursiva genérica com subitens.
  \begin{questoes}
    \item (1,5 ponto) Enunciado do primeiro subitem.
    \item (1,5 ponto) Enunciado do segundo subitem.
  \end{questoes}
  \vspace{4cm} % Espaço para resolução

  % Questão com Fórmula em Bloco
  \item (2,0 pontos) Enunciado de uma questão discursiva genérica contendo uma fórmula em destaque:
  \[ f(x) = \int_{a}^{b} g(x) \, dx \]
  \vspace{5cm}

  % Questão de Verdadeiro ou Falso
  \item (2,5 pontos) Julgue as afirmações genéricas abaixo como Verdadeiras (V) ou Falsas (F):
  \begin{itemize}[label=\(\square\)]
    \item Primeira afirmação para análise.
    \item Segunda afirmação para análise.
    \item Terceira afirmação para análise.
  \end{itemize}

\end{questoes}

\end{document}